% A course handout: learning goals, colored boxes, a worked example,
% a theorem with proof, exercises and a summary.
% Edit the course name and the metadata below.
\documentclass[11pt,parskip=half]{scrartcl}

\usepackage[T1]{fontenc}
\usepackage{lmodern}
\usepackage{microtype}
\usepackage[margin=2.2cm]{geometry}
\usepackage{amsmath,amssymb,amsthm}
\usepackage{enumitem}
\usepackage{xcolor}
\usepackage[most]{tcolorbox}
\usepackage{scrlayer-scrpage}
\usepackage{hyperref}

\definecolor{navy}{HTML}{1F4E79}
\definecolor{forest}{HTML}{2E6B3A}
\definecolor{amber}{HTML}{8A5A00}

% Colored boxes. Use them as \begin{definitionbox}{Name} ... \end{definitionbox}.
\newtcolorbox{definitionbox}[1]{
  colback=navy!5, colframe=navy, fonttitle=\bfseries,
  title={Definition: #1}, breakable
}
\newtcolorbox{examplebox}[1]{
  colback=forest!5, colframe=forest, fonttitle=\bfseries,
  title={Example: #1}, breakable
}
\newtcolorbox{summarybox}{
  colback=amber!6, colframe=amber, fonttitle=\bfseries,
  title={Summary}, breakable
}

\theoremstyle{plain}
\newtheorem{theorem}{Theorem}[section]

% Running header with the course name and the page number.
\pagestyle{scrheadings}
\clearpairofpagestyles
\ihead{Course Name, Term Year}
\ohead{Page \pagemark}

\title{{{title}}}
\subtitle{Lecture handout}
\author{{{author}}}
\date{\today}

\begin{document}

\maketitle

\section{Learning goals}

By the end of this handout you should be able to:
\begin{itemize}
  \item state the definition of a convergent sequence;
  \item prove that a given sequence converges, using the $\varepsilon$-$N$ definition;
  \item explain why every convergent sequence is bounded.
\end{itemize}

\section{Key definitions}

\begin{definitionbox}{Convergent sequence}
  A sequence $(a_n)$ converges to $L$ when, for every $\varepsilon > 0$,
  there is an $N$ such that
  \[
    |a_n - L| < \varepsilon \quad \text{for all } n \ge N.
  \]
\end{definitionbox}

\begin{definitionbox}{Bounded sequence}
  A sequence $(a_n)$ is bounded when there is a constant $M$ with
  $|a_n| \le M$ for every $n$.
\end{definitionbox}

\section{Worked example}

\begin{examplebox}{A rational sequence}
  Show that $a_n = \dfrac{2n+1}{n+3}$ converges to $2$.

  Compute the distance to the limit:
  \[
    |a_n - 2| = \left|\frac{2n+1 - 2(n+3)}{n+3}\right| = \frac{5}{n+3}.
  \]
  Given $\varepsilon > 0$, the inequality $\frac{5}{n+3} < \varepsilon$ holds
  as soon as $n > \frac{5}{\varepsilon} - 3$. Choose $N = \left\lceil \frac{5}{\varepsilon} \right\rceil$.
  Then $n \ge N$ implies $|a_n - 2| < \varepsilon$, so $a_n \to 2$.
\end{examplebox}

\section{A theorem with proof}

\begin{theorem}[Convergent sequences are bounded]\label{thm:bounded}
  If $(a_n)$ converges, then $(a_n)$ is bounded.
\end{theorem}

\begin{proof}
  Let $L$ be the limit. Taking $\varepsilon = 1$ in the definition gives an
  index $N$ with $|a_n - L| < 1$ for all $n \ge N$. Hence $|a_n| < |L| + 1$
  for those $n$. The finitely many terms $a_1, \dots, a_{N-1}$ are bounded
  by their largest absolute value. Taking $M$ to be the larger of that value
  and $|L| + 1$ bounds the whole sequence.
\end{proof}

\section{Exercises}

\begin{enumerate}
  \item Show that $\dfrac{1}{n}$ converges to $0$.
  \item Show that if $a_n \to L$ and $b_n \to K$, then $a_n + b_n \to L + K$.
  \item Find the limit of $\dfrac{3n^2 + 1}{n^2 + 5}$ and justify your answer.
\end{enumerate}

\begin{summarybox}
  \begin{itemize}
    \item A sequence converges when its terms eventually stay within any
      distance $\varepsilon$ of the limit.
    \item Convergence implies boundedness, but a bounded sequence need not converge.
    \item Every example starts from the definition, so practice with
      $\varepsilon$-$N$ arguments before moving on.
  \end{itemize}
\end{summarybox}

\end{document}
